%============================================================
%  CHAPITRE II — État de l'Art
%  Développement d'un Firmware Générique, Modulaire et
%  Configurable pour Appareils Télématiques
%  Stagiaire : Ihssan OULGHAZI — Smart Automation Technologies
%============================================================

\chapter{État de l'Art et Fondements Techniques }

%---- Sommaire local ----------------------------------------
\section*{Sommaire}
\begin{quote}
\small
\textbf{1} \hspace{1em} Introduction \dotfill \pageref{sec:intro_ch2}\\
\textbf{2} \hspace{1em} Généralités sur les firmwares embarqués \dotfill \pageref{sec:generalites}\\
\hspace{2.5em} 2.1 \hspace{0.5em} Définition d'un firmware \dotfill \pageref{subsec:def_fw}\\
\hspace{2.5em} 2.2 \hspace{0.5em} Rôle du firmware dans un appareil télématique \dotfill \pageref{subsec:role_fw}\\
\hspace{2.5em} 2.3 \hspace{0.5em} Firmware spécifique vs firmware générique \dotfill \pageref{subsec:specifique_vs_generique}\\
\hspace{2.5em} 2.4 \hspace{0.5em} Enjeux de maintenabilité et d'évolutivité \dotfill \pageref{subsec:maintenabilite}\\
\textbf{3} \hspace{1em} Plateforme ESP32-P4 et environnement ESP-IDF \dotfill \pageref{sec:esp32}\\
\hspace{2.5em} 3.1 \hspace{0.5em} Présentation de l'ESP32-P4 \dotfill \pageref{subsec:esp32p4}\\
\hspace{2.5em} 3.2 \hspace{0.5em} Mémoire, flash et partitionnement \dotfill \pageref{subsec:memoire}\\
\hspace{2.5em} 3.3 \hspace{0.5em} NVS pour la configuration persistante \dotfill \pageref{subsec:nvs}\\
\hspace{2.5em} 3.4 \hspace{0.5em} ESP-IDF comme framework de développement \dotfill \pageref{subsec:espidf}\\
\hspace{2.5em} 3.5 \hspace{0.5em} Chaîne de compilation, structure et conventions \dotfill \pageref{subsec:toolchain}\\
\textbf{4} \hspace{1em} Systèmes temps réel et FreeRTOS \dotfill \pageref{sec:freertos}\\
\hspace{2.5em} 4.1 \hspace{0.5em} Notion de RTOS \dotfill \pageref{subsec:rtos}\\
\hspace{2.5em} 4.2 \hspace{0.5em} Tâches, priorités et ordonnancement \dotfill \pageref{subsec:tasks}\\
\hspace{2.5em} 4.3 \hspace{0.5em} Queues, timers et event groups \dotfill \pageref{subsec:sync}\\
\hspace{2.5em} 4.4 \hspace{0.5em} Watchdog et supervision interne \dotfill \pageref{subsec:watchdog}\\
\hspace{2.5em} 4.5 \hspace{0.5em} Contraintes mémoire et stack \dotfill \pageref{subsec:mem_stack}\\
\textbf{5} \hspace{1em} Architectures logicielles embarquées modulaires \dotfill \pageref{sec:modularity}\\
\hspace{2.5em} 5.1 \hspace{0.5em} Architecture en couches \dotfill \pageref{subsec:layers}\\
\hspace{2.5em} 5.2 \hspace{0.5em} Séparation driver / service \dotfill \pageref{subsec:driver_service}\\
\hspace{2.5em} 5.3 \hspace{0.5em} HAL — Hardware Abstraction Layer \dotfill \pageref{subsec:hal}\\
\hspace{2.5em} 5.4 \hspace{0.5em} Interfaces de modules : init, start, stop, handle\_event \dotfill \pageref{subsec:interfaces}\\
\hspace{2.5em} 5.5 \hspace{0.5em} Patterns event-driven et pub/sub \dotfill \pageref{subsec:eventdriven}\\
\hspace{2.5em} 5.6 \hspace{0.5em} Pattern strategy pour la sélection runtime \dotfill \pageref{subsec:strategy}\\
\textbf{6} \hspace{1em} Services fondamentaux d'un firmware télématique \dotfill \pageref{sec:services}\\
\hspace{2.5em} 6.1 \hspace{0.5em} Gestion de la configuration \dotfill \pageref{subsec:config}\\
\hspace{2.5em} 6.2 \hspace{0.5em} Stockage local et endurance flash \dotfill \pageref{subsec:stockage}\\
\hspace{2.5em} 6.3 \hspace{0.5em} Ring buffer pour logs et données offline \dotfill \pageref{subsec:ringbuffer}\\
\hspace{2.5em} 6.4 \hspace{0.5em} Logging structuré \dotfill \pageref{subsec:logging}\\
\hspace{2.5em} 6.5 \hspace{0.5em} Power management \dotfill \pageref{subsec:power}\\
\hspace{2.5em} 6.6 \hspace{0.5em} Diagnostic et télémétrie interne \dotfill \pageref{subsec:diag}\\
\hspace{2.5em} 6.7 \hspace{0.5em} Gestion des erreurs et politiques d'incidents \dotfill \pageref{subsec:erreurs}\\
\textbf{7} \hspace{1em} Communication et configurabilité protocolaire \dotfill \pageref{sec:comms}\\
\hspace{2.5em} 7.1 \hspace{0.5em} Besoins de communication dans un appareil télématique \dotfill \pageref{subsec:besoin_comms}\\
\hspace{2.5em} 7.2 \hspace{0.5em} Protocoles TCP, UDP et MQTT \dotfill \pageref{subsec:protocoles}\\
\hspace{2.5em} 7.3 \hspace{0.5em} Interface CommsProvider \dotfill \pageref{subsec:commsprovider}\\
\hspace{2.5em} 7.4 \hspace{0.5em} Sélection du protocole par configuration \dotfill \pageref{subsec:proto_select}\\
\textbf{8} \hspace{1em} Tests et durcissement des firmwares embarqués \dotfill \pageref{sec:tests}\\
\hspace{2.5em} 8.1 \hspace{0.5em} Tests unitaires \dotfill \pageref{subsec:unit_tests}\\
\hspace{2.5em} 8.2 \hspace{0.5em} Tests d'intégration \dotfill \pageref{subsec:integration_tests}\\
\hspace{2.5em} 8.3 \hspace{0.5em} Tests de reboot et de configuration corrompue \dotfill \pageref{subsec:reboot_tests}\\
\hspace{2.5em} 8.4 \hspace{0.5em} Profiling RAM/stack \dotfill \pageref{subsec:profiling}\\
\hspace{2.5em} 8.5 \hspace{0.5em} Release candidate et documentation \dotfill \pageref{subsec:release}\\
\textbf{9} \hspace{1em} Conclusion \dotfill \pageref{sec:conclu_ch2}\\
\end{quote}

\newpage

%============================================================
\section{Introduction}
\label{sec:intro_ch2}
%============================================================
Le développement d'un firmware générique, modulaire et configurable pour appareils télématiques embarqués mobilise plusieurs disciplines : systèmes temps réel, architecture logicielle embarquée, gestion de la persistance et protocoles IoT. Ce chapitre pose l'état des connaissances nécessaire pour comprendre les choix techniques du projet, en suivant l'architecture du firmware développé : notions fondamentales de firmware et enjeux de maintenabilité, plateforme ESP32-P4 et environnement ESP-IDF, programmation temps réel avec FreeRTOS, architecture logicielle modulaire (HAL et patrons de conception adaptés), les six services fondamentaux d'un firmware télématique, la communication protocolaire, et enfin les stratégies de test et de durcissement, avant une conclusion synthétisant les choix technologiques retenus.

%============================================================
\section{Généralités sur les firmwares embarqués}
\label{sec:generalites}
%============================================================

%------------------------------------------------------------
\subsection{Définition d'un firmware}
\label{subsec:def_fw}
%------------------------------------------------------------

Le terme \textit{firmware} désigne le logiciel de bas niveau qui réside en
mémoire non-volatile au sein d'un système électronique et qui prend en charge
le contrôle direct du matériel sous-jacent. Contrairement aux logiciels
applicatifs qui s'exécutent sur un système d'exploitation généraliste avec
des couches d'abstraction multiples, le firmware est intrinsèquement lié à
la plateforme matérielle sur laquelle il s'exécute : il initialise les
périphériques, configure les horloges système, orchestre les interruptions
et expose des services de plus haut niveau aux composants logiciels qui
s'appuient sur lui.

Dans les systèmes embarqués modernes, le firmware occupe une position
architecturale intermédiaire entre le matériel brut et la logique applicative
métier. Il est stocké dans une mémoire Flash interne ou externe, chargé au
démarrage par un bootloader, et s'exécute directement depuis cette mémoire
ou après copie partielle en RAM selon les besoins de performance. Sa mise à
jour, historiquement délicate car nécessitant un accès physique à l'appareil,
est aujourd'hui facilitée par les mécanismes de mise à jour Over-The-Air
(OTA) qui permettent de déployer une nouvelle image firmware à distance, via
le réseau cellulaire ou Wi-Fi.

Un firmware embarqué bien conçu se distingue par trois propriétés
fondamentales. La \textbf{fiabilité} est la première : le firmware doit
fonctionner sans interruption non planifiée dans des conditions
environnementales parfois sévères (vibrations, écarts de température,
coupures d'alimentation brusques). La \textbf{déterminisme temporel} est
la deuxième : les opérations critiques doivent s'exécuter dans des délais
bornés et prévisibles, ce qui distingue le firmware des logiciels de bureau
où une latence occasionnelle est acceptable. Enfin, la \textbf{frugalité
en ressources} est la troisième : les microcontrôleurs embarqués disposent
de quelques centaines de kilo-octets de RAM et d'un ou deux mégaoctets de
Flash, imposant une gestion rigoureuse de chaque ressource allouée.

%------------------------------------------------------------
\subsection{Rôle du firmware dans un appareil télématique}
\label{subsec:role_fw}

Au sein d'un dispositif télématique embarqué dans un véhicule, le firmware
assure un ensemble de tâches qui couvrent l'ensemble du cycle de vie de
l'information, depuis sa captation jusqu'à son acheminement vers
l'infrastructure distante.

La première mission concerne la \textbf{collecte des données}. Le firmware
interagit avec le module GNSS afin de récupérer les informations de
position, sollicite les capteurs inertiels pour repérer les comportements de
conduite (accélérations soudaines, freinages appuyés, virages serrés), et
exploite les liaisons CAN Bus ou OBD-II pour relever des paramètres tels que
le régime moteur, la vitesse, la consommation de carburant ou les codes
d'anomalie.

La deuxième mission porte sur le \textbf{prétraitement embarqué}. Avant toute
communication vers l'extérieur, le firmware met en {\oe}uvre des opérations
de nettoyage, d'identification d'événements et de regroupement des mesures
brutes. Il établit une distinction entre les données nécessitant un envoi
instantané, comme les alertes urgentes, et celles pouvant être consolidées
dans des envois espacés afin de réduire la consommation en données mobiles.

La troisième mission concerne le \textbf{pilotage de la liaison réseau}. Le
firmware initie et maintient la connexion au réseau cellulaire, ouvre des
sessions de communication via MQTT ou TCP vers les serveurs de la
plateforme, et assure la reprise automatique en cas de coupure du signal.
Une mémoire tampon locale conserve par ailleurs les données collectées
durant les interruptions de réseau, évitant ainsi toute perte
d'information.

La quatrième mission, enfin, relève de la \textbf{surveillance interne du
dispositif}. Le firmware contrôle en continu l'état de ses différents
sous-systèmes : occupation de la pile des tâches, cohérence des paramètres
enregistrés en mémoire non volatile, niveau de la batterie de secours, et
fonctionnement des périphériques connectés. Lorsqu'une défaillance est
repérée, des mécanismes de récupération adaptés sont déclenchés avant
notification de la plateforme distante.
\begin{figure}[H]
    \centering
    \includegraphics[width=0.90\textwidth]{figures/fw_responsibilities.png}
    \caption{Responsabilités fonctionnelles d'un firmware télématique embarqué}
    \label{fig:fw_responsibilities}
\end{figure}

L'articulation de ces quatre missions --- collecte des données,
prétraitement embarqué, pilotage de la liaison réseau et surveillance
interne --- ne constitue pas un simple inventaire de fonctions isolées,
mais reflète l'organisation interne même du firmware développé. Comme le
montre la Fig.~\ref{fig:fw_responsibilities}, chacune de ces responsabilités
correspond à un bloc fonctionnel distinct gravitant autour du module central
\emph{ESP32-P4 Firmware} : le bloc \emph{Acquisition et traitement} prend en
charge la captation et le prétraitement des données issues du GNSS, de
l'IMU et du bus véhicule ; le bloc \emph{Connectivité et sécurité} assure le
pilotage de la liaison réseau via GSM/GPRS et MQTT ; le bloc \emph{Gestion
système et services} implémente les mécanismes de surveillance interne
(watchdog, power manager, health counters) ; et le bloc \emph{Logique de
décision et alertes} exploite les données traitées pour générer les
notifications transmises vers le serveur de fleet management. Cette
correspondance directe entre les responsabilités fonctionnelles décrites et
l'architecture représentée sur la figure illustre concrètement le rôle
central que joue le firmware dans la chaîne de valeur d'un appareil
télématique embarqué.



%------------------------------------------------------------
\subsection{Firmware spécifique versus firmware générique}
\label{subsec:specifique_vs_generique}
%------------------------------------------------------------

L'approche traditionnelle de développement firmware consiste à produire un
logiciel étroitement couplé à une configuration matérielle précise et à
un cas d'usage déterminé. Cette approche dite \textbf{spécifique} présente
l'avantage d'une optimisation maximale pour le scénario cible : le code
est minimal, sans abstractions inutiles, et chaque ressource matérielle
est exploitée au maximum de ses capacités. Elle est adaptée aux produits
à volume élevé dont la configuration matérielle ne variera jamais.

Cependant, pour un fabricant qui propose plusieurs variantes d'un même
appareil télématique --- différant par leur module de connectivité, leur
interface bus véhicule ou leur dotation en capteurs --- l'approche spécifique
engendre une multiplication des bases de code. Chaque variante possède son
propre firmware, développé et maintenu indépendamment. Les corrections de
bogues doivent être reportées manuellement sur chaque variante, les nouvelles
fonctionnalités doivent être re-développées autant de fois qu'il existe de
variantes, et les tests de validation doivent être répétés sur l'ensemble
du parc matériel. Le coût de maintenance croît linéairement avec le nombre
de variantes, rendant rapidement l'approche ingérable.

L'\textbf{approche générique} répond à cette problématique en concevant
le firmware comme une plateforme logicielle commune, configurable à la
compilation et à l'exécution pour s'adapter aux différentes variantes
matérielles. Le principe central est de maximiser la réutilisation du code
entre les variantes, en isolant les parties dépendantes du matériel dans
des modules d'adaptation (la HAL) et en paramétrant les comportements
variables via un système de configuration. Une seule base de code est ainsi
maintenue, et les corrections bénéficient immédiatement à toutes les
variantes.

La contrepartie de l'approche générique est une complexité architecturale
plus élevée. La conception des abstractions, des interfaces inter-modules
et du système de configuration demande un investissement initial significatif.
De plus, les abstractions introduisent un overhead en mémoire et en temps
d'exécution, qui doit être évalué et maintenu dans des limites acceptables
pour les microcontrôleurs contraints ciblés.
\newpage
La figure~\ref{fig:specific_vs_generic} illustre cette différence
architecturale : à gauche, l'approche spécifique avec ses bases de code
multiples liées chacune à un matériel précis ; à droite, l'approche
générique fondée sur une base de code commune, une couche de configuration
et une couche d'abstraction matérielle (HAL) permettant de dériver les
différentes variantes du produit.

\begin{figure}[H]
    \centering
    \includegraphics[width=0.80\textwidth]{figures/specific_vs_generic.png}
    \caption{Comparaison entre approche firmware spécifique et approche générique}
    \label{fig:specific_vs_generic}
\end{figure}



%------------------------------------------------------------
\subsection{Enjeux de maintenabilité et d'évolutivité}
\label{subsec:maintenabilite}
%------------------------------------------------------------

La maintenabilité et l'évolutivité d'un firmware embarqué sont des propriétés
qui conditionnent directement le coût total de possession d'un produit sur
son cycle de vie. Un firmware difficile à maintenir accumule progressivement
une dette technique qui ralentit chaque évolution et augmente le risque
d'introduction de régressions. Dans le domaine des appareils télématiques,
où les déploiements sur le terrain peuvent atteindre plusieurs milliers
d'unités et s'étendre sur cinq à dix ans, ces propriétés revêtent une
importance stratégique.

La \textbf{maintenabilité} d'un firmware est favorisée par plusieurs
pratiques architecturales. La séparation des préoccupations (\textit garantit que chaque module a une responsabilité unique et bien
délimitée, ce qui réduit la surface d'impact de toute modification. Les
interfaces contractuelles entre modules permettent de remplacer une
implémentation par une autre sans affecter les modules clients. La
traçabilité des événements par un système de journalisation structuré
facilite le diagnostic des comportements anormaux sur le terrain.

L'\textbf{évolutivité} désigne la capacité du firmware à intégrer de
nouvelles fonctionnalités ou à supporter de nouveaux matériels sans
nécessiter une refonte architecturale. Elle est garantie par la conception
d'interfaces d'extension ouvertes, par le recours à des mécanismes de
découverte dynamique des modules présents (registre de capabilities) et
par une gestion de la configuration qui accepte l'ajout de nouveaux
paramètres sans invalider les configurations existantes.
\newpage
%============================================================
\section{Plateforme ESP32-P4 et environnement ESP-IDF}
\label{sec:esp32}
%============================================================


\subsection{Présentation de l'ESP32-P4 et justification du choix}
\label{subsec:esp32p4}

L'ESP32-P4 est un System-on-Chip (SoC) haute performance développé
par Espressif Systems, conçu pour les applications IoT embarquées
nécessitant une capacité de traitement élevée combinée à une gestion
fine de l'énergie \cite{espressif_esp32p4}. Son choix dans le cadre
du projet \textit{GlobalFleet} a été guidé par une analyse comparative
avec les alternatives du marché, résumée dans le
tableau~\ref{tab:soc_comparison}.

\begin{table}[H]
\centering
\renewcommand{\arraystretch}{1.6}
\caption{Comparaison des plateformes SoC candidates pour le projet \textit{GlobalFleet}}
\label{tab:soc_comparison}
\footnotesize
\begin{tabular}{|p{3.2cm}|p{2.2cm}|p{2.2cm}|p{2.2cm}|p{2.4cm}|}
\hline
\rowcolor{gray!20}
\textbf{Critère} & \textbf{ESP32-P4} & \textbf{STM32H7} & \textbf{ESP32-S3} & \textbf{Nordic nRF9160} \\
\hline
Cœurs processeur & \mbox{2×RISC-V} 400\,MHz & 1×Cortex-M7 480\,MHz & 2×Xtensa LX7 240\,MHz & 1×Cortex-M33 \\
\hline
SRAM interne & 768\,Ko & 1\,Mo & 512\,Ko & 256\,Ko \\
\hline
Coprocesseur LP & Oui (40\,MHz) & Non & Non & Non \\
\hline
Flash Encryption & Oui & Non natif & Oui & Non \\
\hline
Support OTA natif & Oui (dual-bank) & Manuel & Oui & Limité \\
\hline
Interfaces CAN & Oui (TWAI) & Oui (bxCAN) & Non & Non \\
\hline
Écosystème SDK & ESP-IDF (officiel) & STM32CubeIDE & ESP-IDF & nRF SDK \\
\hline
Coût unitaire estimé & Faible & Moyen-élevé & Faible & Élevé \\
\hline
\rowcolor{green!15}
\textbf{Verdict} & \textbf{Retenu} & Non retenu & Écarté & Non retenu \\
\hline
\end{tabular}
\end{table}

L'ESP32-P4 a été retenu pour cinq raisons déterminantes. Premièrement,
son \textbf{coprocesseur LP} (Low Power, 40\,MHz) permet de maintenir
la surveillance des sources de réveil à une consommation inférieure
à 5\,mA sans activer les cœurs RISC-V principaux, ce qui est critique
pour les appareils alimentés sur batterie de secours. Deuxièmement,
son \textbf{contrôleur TWAI natif} (compatible CAN 2.0B et CAN-FD)
élimine la nécessité d'un transceiver externe pour l'interface J1939
véhicule. Troisièmement, le mécanisme \textbf{OTA Dual-Bank} intégré
dans l'ESP-IDF permet des mises à jour atomiques avec rollback
automatique sans développement additionnel. Quatrièmement, la
\textbf{Flash Encryption} et le \textbf{Secure Boot v2} répondent
aux exigences de sécurité croissantes du secteur. Cinquièmement,
l'\textbf{écosystème ESP-IDF mature} avec FreeRTOS SMP natif et une
communauté active réduit le risque projet.
\newpage
Le STM32H7 a été écarté malgré ses performances brutes supérieures
en raison de l'absence de coprocesseur LP natif et de l'écosystème
OTA moins intégré. Le nRF9160, bien que performant pour les
communications LTE-M, présente un coût unitaire prohibitif et ne
supporte pas le CAN Bus natif.

Au-delà des critères purement matériels résumés dans le tableau
\ref{tab:soc_comparison}, le choix définitif de l'ESP32-P4 repose
également sur la maturité et la cohérence de son environnement
logiciel. Le framework ESP-IDF, construit autour de FreeRTOS avec
support multi-cœur natif (SMP), fournit une couche d'abstraction
matérielle homogène à travers les pilotes \textit{esp\_driver}, ce
qui simplifie considérablement le portage du firmware entre
différentes variantes de cartes basées sur la même famille de SoC.
Cette abstraction concerne aussi bien les périphériques de
communication (UART, SPI, I\textsuperscript{2}C, TWAI) que les
modules de gestion d'énergie, ce qui réduit le volume de code
spécifique à chaque révision matérielle.

Par ailleurs, l'intégration native d'un gestionnaire de partitions
Flash permet de séparer proprement le code applicatif, la
configuration persistante stockée dans la partition NVS, et le
système de fichiers embarqué basé sur LittleFS. Cette séparation
facilite la maintenance à distance du parc de boîtiers et limite les
risques de corruption de données lors des mises à jour OTA.
L'écosystème ESP-IDF offre également des outils intégrés de
diagnostic et de débogage (monitor série, core dump, analyse de
piles), ce qui réduit le temps de mise au point lors des phases de
développement et de validation sur le terrain.

Enfin, sur le plan organisationnel, la disponibilité d'une
documentation officielle complète, d'exemples de code maintenus par
Espressif et d'une communauté open-source active constitue un
facteur de réduction du risque projet non négligeable. Pour un
projet comme \textit{GlobalFleet}, où plusieurs variantes matérielles
(modules GPS, modems cellulaires, capteurs CAN, interfaces
d'affichage) doivent pouvoir partager une base logicielle commune
tout en conservant la possibilité de mises à jour différenciées par
configuration, cette homogénéité de plateforme et d'outillage
représente un avantage déterminant par rapport aux alternatives
envisagées.
\begin{figure}[H]
\centering
\includegraphics[width=0.50\textwidth]{figures/esp32p4.png}
\caption{Microcontrôleur ESP32-P4 — plateforme cible retenue pour le firmware
\textit{GlobalFleet}.
\textit{Source : \cite{espressif_esp32p4}}}
\label{fig:esp32p4_photo}
\end{figure}
Sur le plan architectural, l'ESP32-P4 intègre deux cœurs RISC-V
haute performance (HP) cadencés jusqu'à 400\,MHz, accompagnés du
coprocesseur LP fonctionnant indépendamment à 40\,MHz. La mémoire
interne se compose de 768\,Ko de SRAM répartis entre les domaines HP
et LP. La richesse des interfaces — UART, SPI, I\textsuperscript{2}C,
CAN (TWAI), Ethernet MAC, USB-OTG — permet de connecter l'ensemble
des périphériques d'un appareil télématique sans composant d'extension.
\begin{figure}[H]
\centering
\includegraphics[width=0.82\textwidth]{figures/esp32p4_block.png}
\caption{Diagramme bloc de l'ESP32-P4 et organisation de ses domaines d'alimentation.
\textit{Source : \cite{espressif_esp32p4}}}
\label{fig:esp32p4_block}
\end{figure}

\subsection{Mémoire, flash et partitionnement}
\label{subsec:memoire}

L'organisation de la mémoire Flash d'un système ESP32 est un aspect
architectural fondamental qui conditionne la robustesse du firmware.
ESP-IDF utilise un schéma de partitionnement flexible défini par une
table stockée à une adresse fixe dans la Flash. Cette table décrit
l'ensemble des régions et leur affectation, permettant de coexister
plusieurs images firmware, une zone de configuration et des systèmes
de fichiers sur le même composant Flash.

Les partitions classiques d'un firmware ESP32 comprennent : une
partition de bootloader à l'adresse 0x0000, la table de partitions
à 0x8000, deux partitions OTA symétriques (OTA\_0
 \newpage et OTA\_1), une
partition NVS pour la configuration persistante, une partition SPIFFS
ou LittleFS pour le système de fichiers embarqué, et optionnellement
une partition factory pour les paramètres immuables de fabrication.

\begin{figure}[H]
\centering
\includegraphics[width=0.70\textwidth]{figures/flash_partitions.png}
\caption{Organisation typique des partitions Flash d'un firmware ESP32.
\textit{Source : \cite{espidf_nvs}}}
\label{fig:flash_partitions}
\end{figure}

Le mécanisme OTA repose sur le principe de la double banque : pendant
que le firmware actuel s'exécute depuis OTA\_0, le nouveau firmware
est téléchargé et écrit dans OTA\_1. Lorsque l'écriture est validée
par vérification de signature, le bootloader est instructé de démarrer
sur OTA\_1 au prochain redémarrage. En cas d'échec, le bootloader
revient automatiquement sur la partition précédente.

\subsection{NVS pour la configuration persistante}
\label{subsec:nvs}

Le système \textbf{NVS} (\textit{Non-Volatile Storage}) d'ESP-IDF est
une bibliothèque de stockage clé-valeur conçue pour persister des
paramètres de configuration dans la Flash interne \cite{espidf_nvs}.
Elle répond à deux contraintes majeures : la durabilité face aux
coupures d'alimentation et la longévité du composant Flash malgré
des mises à jour fréquentes.

Le NVS organise les données en \textit{namespaces}, espaces de noms
logiques permettant à des modules indépendants de gérer leur propre
espace sans risque de collision de clés. Chaque entrée est une paire
clé-valeur où la clé est une chaîne de 15 caractères maximum et la
valeur peut être un entier de 8 à 64 bits, une chaîne ou un blob
binaire permettant de stocker des structures C complètes.

Le NVS implémente un algorithme de \textit{wear leveling} par rotation
des pages : chaque page est utilisée séquentiellement, et lorsqu'une
page est pleine, ses entrées valides sont compactées vers une nouvelle
page avant effacement de l'ancienne. La résistance aux coupures
d'alimentation est assurée par l'écriture atomique de chaque entrée.

\subsection{ESP-IDF comme framework de développement}
\label{subsec:espidf}

L'\textbf{ESP-IDF} (\textit{Espressif IoT Development Framework}) est
le SDK officiel d'Espressif, qui constitue bien plus qu'un ensemble
de bibliothèques périphériques \cite{espidf_github} : c'est un
environnement complet intégrant le noyau FreeRTOS, une couche HAL
structurée, un système de composants réutilisables, un système de
configuration Kconfig, et les outils de compilation, flashage et
débogage.

L'architecture d'ESP-IDF est organisée en composants (\textit{components}),
chacun constituant un module logiciel autonome avec ses propres
sources, ses en-têtes publics et son fichier Kconfig. Les composants
peuvent déclarer des dépendances explicites, ce qui permet au système
de build CMake de construire automatiquement le graphe de dépendances
et d'assurer la compilation dans le bon ordre. Cette organisation
est directement réutilisée dans le firmware \textit{GlobalFleet}.

\begin{figure}[H]
\centering
\includegraphics[width=0.60\textwidth]{figures/espidf_archi.png}
\caption{Architecture en composants d'ESP-IDF et intégration de FreeRTOS.
\textit{Source : \cite{espidf_github}}}
\label{fig:espidf_archi}
\end{figure}
\newpage
ESP-IDF intègre une version enrichie de FreeRTOS supportant le mode
SMP (\textit{Symmetric Multi-Processing}) sur les ESP32 à double cœur,
permettant d'épingler une tâche sur un cœur spécifique via le paramètre
\texttt{xCoreID} de \texttt{xTaskCreatePinnedToCore()}. Cette capacité
est exploitée dans le firmware \textit{GlobalFleet} pour séparer les
tâches de communication réseau (cœur 0) des tâches d'acquisition de
données (cœur 1).

\subsection{Chaîne de compilation, structure de projet et conventions}
\label{subsec:toolchain}

La chaîne de compilation d'ESP-IDF repose sur CMake comme système
de build et sur une toolchain GCC RISC-V croisée fournie par Espressif.
La structure d'un projet ESP-IDF suit des conventions précises : le
répertoire racine contient \texttt{CMakeLists.txt}, le fichier
\texttt{sdkconfig} généré par menuconfig, et le répertoire
\texttt{main/} contenant le point d'entrée. Les composants
additionnels résident dans \texttt{components/}.

Le système \textbf{Kconfig} permet de définir des options de
configuration hiérarchiques exposées dans le menu \texttt{idf.py
menuconfig}. Ces options génèrent des macros préprocesseur dans
\texttt{sdkconfig.h}, permettant d'inclure ou d'exclure
conditionnellement des blocs de code à la compilation selon la
configuration cible. Cette mécanique est au cœur du mécanisme de
configurabilité à la compilation du firmware \textit{GlobalFleet}.

%============================================================
\section{Systèmes temps réel et FreeRTOS}
\label{sec:freertos}
%============================================================

\subsection{Notion de RTOS et justification du choix}
\label{subsec:rtos}

Un système d'exploitation temps réel (RTOS) est un noyau logiciel
conçu pour garantir que les contraintes temporelles d'un système
embarqué sont respectées de manière déterministe et prévisible.
La distinction fondamentale avec un OS généraliste ne réside pas dans
la vitesse d'exécution, mais dans la \textbf{prévisibilité} : un RTOS
garantit qu'une tâche de haute priorité sera activée dans un délai
maximal borné après l'événement déclencheur \cite{freertos_official}.

Cette garantie est réalisée grâce à un ordonnanceur préemptif à
priorités fixes : dès qu'une tâche de priorité supérieure devient
prête, l'ordonnanceur interrompt immédiatement la tâche en cours.
On distingue le temps réel strict (\textit{hard real-time}), où tout
dépassement de l'échéance est une défaillance inacceptable, du temps
réel souple (\textit{soft real-time}), où des dépassements
occasionnels sont tolérés. Dans le firmware \textit{GlobalFleet},
la majorité des contraintes relèvent du temps réel souple.
\newpage
Le tableau~\ref{tab:rtos_comparison} présente la comparaison des
RTOS candidats qui ont été évalués avant de retenir FreeRTOS.

\begin{table}[H]
\centering
\renewcommand{\arraystretch}{1.6}
\caption{Comparaison des RTOS candidats pour le projet \textit{GlobalFleet}}
\label{tab:rtos_comparison}
\footnotesize
\begin{tabular}{|p{3.0cm}|p{2.3cm}|p{2.3cm}|p{2.3cm}|p{2.3cm}|}
\hline
\rowcolor{gray!20}
\textbf{Critère} & \textbf{FreeRTOS} & \textbf{Zephyr} & \textbf{RTX (Azure)} & \textbf{ESP-IDF bare} \\
\hline
Intégration ESP-IDF & Native (officielle) & Partielle & Non & N/A \\
\hline
Support SMP dual-core & Oui (v10.5) & Oui & Non & Non \\
\hline
Empreinte mémoire & Très faible & Moyenne & Faible & Minimale \\
\hline
Licence & MIT & Apache 2.0 & MIT & — \\
\hline
Maturité industrielle & Très élevée & Élevée & Moyenne & — \\
\hline
Documentation & Complète \cite{freertos_book} & Complète & Limitée & — \\
\hline
\rowcolor{green!15}
\textbf{Verdict} & \textbf{Retenu} & Non retenu & Non retenu & Non retenu \\
\hline
\end{tabular}
\end{table}

FreeRTOS a été retenu pour quatre raisons. Premièrement, son
\textbf{intégration native dans ESP-IDF} élimine toute problématique
de portage : l'environnement de développement Espressif livre FreeRTOS
v10.5.1 préconfiguré et optimisé pour l'ESP32-P4. Deuxièmement, son
\textbf{support SMP natif} permet d'exploiter les deux cœurs RISC-V
de l'ESP32-P4 sans développement additionnel. Troisièmement, sa
\textbf{très faible empreinte mémoire} (quelques dizaines de Ko)
laisse le maximum de SRAM disponible pour les buffers applicatifs.
Quatrièmement, sa \textbf{maturité industrielle} et sa documentation
exhaustive réduisent le risque de comportements non documentés en
production \cite{freertos_official}.

\begin{figure}[H]
\centering
\includegraphics[width=0.90\textwidth]{figures/rtos_preemption.png}
\caption{Mécanisme de préemption dans un RTOS à priorités fixes.
\textit{Source : adapté de \cite{freertos_book}}}
\label{fig:rtos_preemption}
\end{figure}

\subsection{Tâches, priorités et ordonnancement}
\label{subsec:tasks}

La tâche (\textit{task}) est l'unité d'exécution fondamentale de
FreeRTOS. Chaque tâche est une fonction C s'exécutant dans une boucle
infinie, dotée de son propre contexte d'exécution comprenant une pile
dédiée, une copie des registres du processeur et un bloc de contrôle
(TCB). L'isolation des contextes garantit qu'une tâche ne peut pas
corrompre accidentellement la pile d'une autre.

Chaque tâche se voit attribuer une \textbf{priorité numérique} au
moment de sa création. L'ordonnanceur sélectionne à chaque instant
la tâche dans l'état \textit{Ready} ayant la priorité la plus haute.
Si plusieurs tâches partagent la même priorité maximale, elles sont
exécutées en tourniquet (\textit{Round-Robin}) par tranche de temps
égale au quantum de l'ordonnanceur (\texttt{configTICK\_RATE\_HZ}
= 1000\,Hz dans ESP-IDF, soit 1\,ms).

Une tâche peut se trouver dans l'un de quatre états :
\textbf{Running} (en cours d'exécution), \textbf{Ready} (prête,
en attente de sélection), \textbf{Blocked} (en attente d'un
événement, sans consommer de CPU) ou \textbf{Suspended} (mise en
pause explicitement). Le passage de Blocked à Ready peut être
déclenché par l'expiration d'un délai, la réception d'un message
ou la libération d'un sémaphore \cite{freertos_book}.

\begin{figure}[H]
\centering
\includegraphics[width=0.90\textwidth]{figures/task_states.png}
\caption{Diagramme d'états d'une tâche FreeRTOS.
\textit{Source : adapté de \cite{freertos_book}}}
\label{fig:task_states}
\end{figure}

\subsection{Queues, timers et event groups}
\label{subsec:sync}

FreeRTOS fournit un ensemble de primitives de synchronisation et de
communication inter-tâches, conçues pour être thread-safe par nature
et supportant l'utilisation depuis les routines d'interruption (ISR)
via des variantes suffixées \texttt{FromISR} \cite{freertos_official}.

La \textbf{file de messages} (\textit{queue}) implémente un tampon
FIFO de capacité fixe, où chaque entrée est une copie de la donnée
envoyée (pas un pointeur), garantissant la sûreté mémoire. Une tâche
tentant de lire dans une file vide se bloque automatiquement pour
le délai spécifié.

Les \textbf{sémaphores} servent à la synchronisation d'événements et
à la protection de ressources partagées. Le \textbf{mutex} est une
variante augmentée d'un mécanisme d'héritage de priorité résolvant
le problème classique d'inversion de priorité.

Les \textbf{groupes d'événements} (\textit{Event Groups}) permettent
à une tâche d'attendre simultanément plusieurs conditions logiques
exprimées comme un masque de bits (opérateur ET ou OU logique). Dans
le firmware \textit{GlobalFleet}, ils implémentent des barrières de
synchronisation au démarrage.

Les \textbf{timers logiciels} permettent d'exécuter des fonctions
de rappel après un délai ou de façon périodique. Ils gèrent dans le
firmware les envois périodiques de heartbeat et les rafraîchissements
des compteurs de santé des modules.

\subsection{Watchdog et supervision interne}
\label{subsec:watchdog}

Les mécanismes de watchdog sont essentiels à la robustesse d'un
firmware embarqué déployé sur le terrain, où une intervention humaine
en cas de blocage logiciel n'est pas envisageable. L'ESP32-P4 intègre
deux watchdogs matériels indépendants \cite{espressif_esp32p4}.

Le \textbf{TWDT} (\textit{Task Watchdog Timer}) surveille que chaque
tâche supervisée effectue régulièrement son appel de rafraîchissement
(\textit{kick}), détectant ainsi les blocages au niveau applicatif.
L'\textbf{IWDT} (\textit{Interrupt Watchdog Timer}) surveille que
l'ordonnanceur FreeRTOS n'est pas bloqué dans une interruption trop
longtemps.
\newpage
Au-dessus des watchdogs matériels, le firmware \textit{GlobalFleet}
implémente une supervision logicielle à deux niveaux : un système de
\textbf{health counters} par module et un \textbf{watchdog logiciel
applicatif} multi-tâches dans le composant \texttt{watchdog\_task}.

\subsection{Contraintes mémoire et stack}
\label{subsec:mem_stack}

La gestion de la mémoire dans un firmware FreeRTOS embarqué est une
contrainte permanente. Contrairement aux logiciels s'exécutant sur
des systèmes avec mémoire virtuelle, un firmware embarqué dispose
d'une quantité fixe et limitée de RAM. Un dépassement de pile
(\textit{stack overflow}) ou une fragmentation excessive du tas
peuvent conduire à des comportements indéfinis.

Chaque tâche FreeRTOS se voit allouer une pile de taille fixe au
moment de sa création. La fonction
\texttt{uxTaskGetStackHighWaterMark()} retourne le nombre minimum
de mots de pile jamais utilisés simultanément, permettant d'affiner
les dimensionnements après des tests en conditions réelles.

Dans le firmware \textit{GlobalFleet}, une politique de
dimensionnement conservatrice est adoptée lors du développement, puis
affinée par mesure systématique des watermarks en phase de test. Les
allocations dynamiques de grande taille sont orientées vers la PSRAM
externe via l'attribut \texttt{MALLOC\_CAP\_SPIRAM}, préservant la
SRAM interne pour les allocations latence-critiques.

%============================================================
\section{Architectures logicielles embarquées modulaires}
\label{sec:modularity}
%============================================================

\subsection{Architecture en couches}
\label{subsec:layers}

L'architecture logicielle en couches est le patron organisationnel
le plus fondamental du génie logiciel embarqué. Son principe consiste
à décomposer le logiciel en niveaux d'abstraction hiérarchiques
superposés, chaque couche n'interagissant qu'avec les couches
immédiatement adjacentes via des interfaces contractuelles définies.
Cette règle d'interaction restreinte confère à l'architecture ses
propriétés essentielles : portabilité, testabilité et maintenabilité.

Dans le firmware \textit{GlobalFleet}, on identifie quatre couches
superposées. La couche \textbf{matérielle} (BSP) expose les primitives
bas niveau d'ESP-IDF et FreeRTOS. La couche \textbf{HAL \& Drivers}
abstrait les périphériques physiques. La couche \textbf{Services Core}
implémente les fonctionnalités transversales (configuration, stockage,
journalisation, communication, gestion d'énergie). La couche
\textbf{Applicative} contient la logique métier spécifique à chaque
type d'appareil.
\newpage
Pour adapter le firmware à une nouvelle variante matérielle, seule
la couche HAL doit être modifiée. Les services core et la logique
applicative restent strictement inchangés.
\begin{figure}[H]
\centering
\includegraphics[width=0.80\textwidth]{figures/layered_archi.png}
\caption{Architecture en couches du firmware générique \textit{GlobalFleet}.
\textit{Source : élaboration propre, inspirée de \cite{espidf_github}}}
\label{fig:layered_archi}
\end{figure}

\noindent
Pour adapter le firmware à une nouvelle variante matérielle, seule
la couche HAL doit être modifiée. Les services core et la logique
applicative restent strictement inchangés.

\subsection{Séparation driver / service}
\label{subsec:driver_service}

La séparation entre le concept de \textbf{driver} et le concept de
\textbf{service} est une distinction architecturale fondamentale dans
un firmware modulaire. Un driver contrôle un périphérique matériel
précis et expose des primitives d'entrée-sortie de bas niveau. Un
service utilise un ou plusieurs drivers pour réaliser une fonction
métier.

Cette séparation présente plusieurs avantages : un driver peut être
remplacé par un stub de test, un même service peut fonctionner avec
différentes implémentations de drivers, et la responsabilité de chaque
composant est clairement délimitée.

\subsection{HAL — Hardware Abstraction Layer}
\label{subsec:hal}

La couche d'abstraction matérielle (HAL) est le mécanisme central
permettant au firmware de fonctionner sur plusieurs variantes
matérielles sans duplication du code. Elle définit des interfaces
abstraites — typiquement des structures de pointeurs de fonctions
en C — que chaque implémentation concrète doit satisfaire.

Le principe de conception de la HAL repose sur l'inversion de
dépendance : les modules de haut niveau dépendent d'abstractions
et non d'implémentations concrètes. Dans le firmware
\textit{GlobalFleet}, la HAL s'appuie sur les drivers ESP-IDF comme
couche intermédiaire entre le matériel brut et les services, avec
une surcouche supplémentaire pour permettre la substitution par des
mocks lors des tests unitaires.

\subsection{Interfaces de modules : init, start, stop, handle\_event}
\label{subsec:interfaces}

Chaque module du firmware \textit{GlobalFleet} implémente un contrat
d'interface standardisé composé de quatre opérations fondamentales.

L'opération \textbf{\texttt{init}} alloue les ressources statiques
du module et enregistre les handlers d'événements auprès du bus
d'événements central. L'opération \textbf{\texttt{start}} crée les
tâches FreeRTOS du module. L'opération \textbf{\texttt{stop}} suspend
proprement le fonctionnement du module sans détruire ses ressources.
L'opération \textbf{\texttt{handle\_event}} reçoit et traite les
événements publiés par d'autres modules via le bus d'événements.

\subsection{Patterns event-driven et pub/sub}
\label{subsec:eventdriven}

L'architecture événementielle (\textit{event-driven architecture})
est un paradigme où les modules ne s'invoquent pas directement mais
communiquent en publiant et en s'abonnant à des événements via un
médiateur central. Ce paradigme est particulièrement bien adapté aux
systèmes embarqués temps réel car il découple naturellement les
producteurs de données de leurs consommateurs.

Dans le firmware \textit{GlobalFleet}, un composant \textbf{Event Bus}
central implémente ce mécanisme. Chaque module peut publier des
événements typés (\texttt{EVENT\_GNSS\_FIX\_ACQUIRED},
\texttt{EVENT\_CAN\_ERROR\_DETECTED}, \texttt{EVENT\_POWER\_BATTERY\_LOW})
et s'abonner aux types qui le concernent, sans que le producteur
ait connaissance de l'identité des consommateurs.

\begin{figure}[h]
\centering
\includegraphics[width=0.60\textwidth]{figures/event_bus.jpg}
\caption{Architecture de l'Event Bus inter-modules du firmware
\textit{GlobalFleet} — patron Publish-Subscribe sur 10 composants.
\textit{Source : élaboration propre}}
\label{fig:event_bus}
\end{figure}

\noindent
L'implémentation de l'Event Bus dans un contexte FreeRTOS s'appuie
sur une file de messages protégée par mutex. Une tâche de dispatching
dédiée lit la file et appelle les callbacks des abonnés, garantissant
que ces callbacks s'exécutent dans le contexte du dispatcher et non
dans celui du producteur, évitant ainsi les problèmes de réentrance.

\subsection{Pattern Strategy pour la sélection runtime}
\label{subsec:strategy}

Le patron de conception Stratégie (\textit{Strategy Pattern}) permet
de définir une famille d'algorithmes interchangeables, d'encapsuler
chacun et de les rendre substituables à l'exécution. En C embarqué,
ce patron est implémenté via des tables de pointeurs de fonctions.

Pour la couche de communication du firmware : selon la configuration
lue au démarrage, le module de communication doit utiliser soit MQTT,
soit TCP, soit UDP. Plutôt que de parsemer le code de branchements
conditionnels, le patron Stratégie encapsule chaque protocole dans
une structure \texttt{comms\_provider\_t} contenant des pointeurs vers
les fonctions \texttt{connect}, \texttt{send}, \texttt{receive} et
\texttt{disconnect}. Au démarrage, la stratégie appropriée est
sélectionnée et injectée dans le module de communication.

%============================================================
\section{Services fondamentaux d'un firmware télématique}
\label{sec:services}
%============================================================

\subsection{Gestion de la configuration}
\label{subsec:config}

Le module de gestion de la configuration constitue l'un des services
les plus transversaux du firmware. Sa responsabilité est de fournir
à tous les autres modules un accès unifié et sécurisé aux paramètres
qui gouvernent leur comportement, en distinguant deux catégories selon
leur nature et leur durée de vie.

La \textbf{configuration factory} regroupe les paramètres définis une
fois pour toutes à la fabrication : numéro de série, identifiant
matériel, version du PCB, capacités hardware activées. Ces paramètres
sont stockés dans une partition Flash protégée et ne doivent jamais
être modifiables sur le terrain.

La \textbf{configuration field} regroupe les paramètres ajustables
lors du déploiement : paramètres réseau, intervalles de reporting,
seuils d'événements, activation de fonctionnalités optionnelles.
Ces paramètres sont stockés en NVS avec vérification d'intégrité
par CRC-32 et numéro de version 16 bits \cite{espidf_nvs}.

\subsection{Stockage local et endurance flash}
\label{subsec:stockage}

Le service de stockage local gère la persistance des données
opérationnelles générées par le firmware. La contrainte d'endurance
de la Flash (100\,000 à 300\,000 cycles d'effacement par secteur)
impose un algorithme de \textit{wear leveling} applicatif distribuant
les écritures sur l'ensemble de la partition allouée.

\subsection{Ring buffer pour logs et données offline}
\label{subsec:ringbuffer}

Un tampon circulaire (\textit{ring buffer}) est une structure de
données particulièrement bien adaptée au stockage temporaire de flux
de données dans les systèmes embarqués à ressources limitées. Son
principe repose sur une zone mémoire de taille fixe gérée avec deux
pointeurs indépendants d'écriture et de lecture.

Dans le firmware \textit{GlobalFleet}, le ring buffer est utilisé à
deux niveaux : le stockage offline des données en l'absence de
connectivité, et le tampon mémoire absorbant les rafales de messages
de log sans bloquer les tâches productrices.

\subsection{Logging structuré}
\label{subsec:logging}

Un système de journalisation robuste est indispensable au diagnostic
d'un firmware déployé sur le terrain. La journalisation structurée
produit des entrées de log formatées selon un schéma défini, incluant
systématiquement un horodatage précis, le niveau de sévérité (DEBUG,
INFO, WARNING, ERROR, CRITICAL), l'identifiant du module émetteur
et le message structuré au format JSON.

\subsection{Power management}
\label{subsec:power}

Le gestionnaire d'énergie du firmware implémente une machine à états
avec plusieurs modes de fonctionnement distincts. Le mode
\textbf{ACTIVE} correspond au fonctionnement nominal complet. Le mode
\textbf{REDUCED} réduit la fréquence d'acquisition. Le mode
\textbf{SLEEP} éteint la connectivité cellulaire et les périphériques
non essentiels. Le mode \textbf{DEEP\_SLEEP} (Ultra-Pro uniquement)
coupe tous les domaines HP pour atteindre une consommation inférieure
à 2\,mA \cite{espressif_esp32p4}.
\begin{figure}[H]
    \centering
    \includegraphics[width=0.60\textwidth]{figures/power_fsm.png}
    \caption{Machine à états du gestionnaire d'énergie du firmware \textit{GlobalFleet}. \textit{Source : élaboration propre}}
    \label{fig:power_fsm}
\end{figure}

\subsection{Diagnostic et télémétrie interne}
\label{subsec:diag}

Le service de diagnostic assure la surveillance en temps réel de
l'état de santé du firmware lui-même. Il collecte périodiquement :
la mémoire heap libre par région, le watermark minimum de pile de
chaque tâche critique, le temps de fonctionnement depuis le dernier
démarrage, la cause du dernier reset, la qualité du signal cellulaire
et la charge CPU moyenne. Ces métriques permettent de détecter une
dégradation progressive avant qu'elle ne devienne critique.

\subsection{Gestion des erreurs et politiques d'incidents}
\label{subsec:erreurs}

La gestion des erreurs dans un firmware destiné à fonctionner sans
supervision humaine doit être à la fois systématique et graduée.
Le firmware \textit{GlobalFleet} définit une hiérarchie de politiques
d'incidents adaptées à la sévérité et à la récupérabilité des erreurs :
journalisation pour les erreurs bénignes, redémarrage du module fautif
avec backoff exponentiel pour les erreurs fonctionnelles, et reset
complet du système pour les erreurs système critiques.

%============================================================
\section{Communication et configurabilité protocolaire}
\label{sec:comms}
%============================================================

\subsection{Besoins de communication dans un appareil télématique}
\label{subsec:besoin_comms}

Les exigences sur le canal de communication sont multiples et parfois
contradictoires. La \textbf{fiabilité de livraison} garantit que les
données critiques sont transmises avec acquittement. L'\textbf{efficacité
protocolaire} minimise les en-têtes et compresse les données pour
réduire la consommation cellulaire. La \textbf{sécurité} protège
les données en transit par chiffrement TLS et authentifie la source
\cite{tls_rfc5246}.

\subsection{Protocoles TCP, UDP et MQTT}
\label{subsec:protocoles}

Le protocole \textbf{TCP} est orienté connexion et garantit la
livraison ordonnée et sans erreur de tous les octets transmis. Son
overhead de connexion le rend moins adapté aux environnements cellulaires
à latence variable.

Le protocole \textbf{UDP} est sans connexion, avec overhead minimal,
adapté aux flux de données périodiques où la perte occasionnelle d'un
paquet est préférable à la latence des retransmissions TCP.

\textbf{MQTT} (\textit{Message Queuing Telemetry Transport}) est un
protocole applicatif publish/subscribe conçu spécifiquement pour les
appareils IoT à ressources contraintes sur des réseaux à faible débit
et haute latence \cite{mqtt_spec}\cite{mqtt_v5}. Il s'appuie sur
TCP et introduit un broker central. Les niveaux QoS (0, 1 et 2)
permettent d'ajuster le compromis entre overhead et garantie de
livraison selon la criticité des données.

\subsection{Interface CommsProvider}
\label{subsec:commsprovider}

L'interface \textbf{CommsProvider} encapsule toutes les opérations
de communication réseau derrière un contrat uniforme, permettant
aux modules consommateurs d'envoyer et de recevoir des données sans
connaître le protocole sous-jacent. L'interface expose un ensemble
minimal de primitives : \texttt{comms\_connect()},
\texttt{comms\_publish(topic, data, len, qos)},
\texttt{comms\_subscribe(topic, callback)},
\texttt{comms\_disconnect()} et \texttt{comms\_get\_status()}.

\subsection{Sélection du protocole par configuration}
\label{subsec:proto_select}

La sélection du protocole de communication effectivement utilisé à
l'exécution est pilotée par un paramètre de la configuration field
de l'appareil. Au démarrage, le module de communication lit ce
paramètre et instancie la structure \texttt{comms\_provider\_t}
correspondant à l'implémentation sélectionnée. Ce mécanisme permet
de déployer le même binaire firmware sur des appareils destinés à
différentes infrastructures serveur, en modifiant uniquement la
configuration sans recompilation.

%============================================================
\section{Outils de développement et environnement Visual Studio Code}
\label{sec:outils}
%============================================================

\subsection{Visual Studio Code et framework ESP-IDF}
\label{subsec:vscode}

L'environnement de développement retenu pour le projet
\textit{GlobalFleet} est \textbf{Visual Studio Code} (VS~Code),
utilisé comme éditeur de code associé au \textbf{framework officiel
ESP-IDF} d'Espressif, exploité directement via son interface en ligne
de commande. Ce choix s'inscrit dans une démarche de productivité et
de cohérence avec les pratiques de l'équipe R\&D de Smart Automation
Technologies, en privilégiant un environnement léger, sans dépendance
à des extensions graphiques additionnelles.

VS~Code offre une interface d'édition performante intégrant un
terminal intégré, dans lequel les commandes \texttt{idf.py build},
\texttt{idf.py flash} et \texttt{idf.py monitor} sont exécutées
directement, sans basculer entre plusieurs fenêtres. Les
fonctionnalités IntelliSense pour C/C++ natives de VS~Code, combinées
à la configuration du \texttt{compile\_commands.json} généré par
ESP-IDF, offrent une complétion de code précise sur l'ensemble des
APIs \texttt{esp\_mqtt}, \texttt{nvs\_flash}, \texttt{esp\_ota\_ops}
et des primitives FreeRTOS.

\begin{figure}[!ht]
\centering

\includegraphics[width=0.80\textwidth]{figures/vscode_espidf.png}
\caption{Environnement Visual Studio Code et terminal ESP-IDF —
interface utilisée pour le développement du firmware \textit{GlobalFleet}.
\textit{Source : capture d'écran du projet}}
\label{fig:vscode_espidf}
\end{figure}

\noindent
Le débogage applicatif a été assuré principalement par
l'instrumentation logicielle (logs \texttt{ESP\_LOGx} et moniteur
série \texttt{idf.py monitor}), permettant le suivi des tâches
FreeRTOS, des transitions d'état et des codes d'erreur retournés par
les APIs ESP-IDF, sans recourir à un débogueur matériel JTAG.

\subsection{Comparaison et justification des outils retenus}
\label{subsec:justif_outils}

Le tableau~\ref{tab:tools_comparison} présente la comparaison des
environnements de développement évalués avant de retenir VS~Code
associé au framework ESP-IDF en ligne de commande.

\begin{table}[H]
\centering
\renewcommand{\arraystretch}{1.6}
\caption{Comparaison des environnements de développement candidats}
\label{tab:tools_comparison}
\footnotesize
\begin{tabular}{|p{3.0cm}|p{2.5cm}|p{2.5cm}|p{2.5cm}|p{2.0cm}|}
\hline
\rowcolor{gray!20}
\textbf{Critère} & \textbf{VS Code + ESP-IDF (CLI)} & \textbf{Eclipse + ESP-IDF} & \textbf{STM32CubeIDE} & \textbf{CLion} \\
\hline
Support ESP32-P4 natif & Officiel & Officiel & Non & Via CMake \\
\hline
Légèreté / performance & Très légère & Lourde & Lourde & Moyenne \\
\hline
Configuration requise & Minimale & Importante & Dédiée STM32 & CMake natif \\
\hline
Intégration Git & Intégrée & Partielle & Partielle & Complète \\
\hline
Débogage & Logs / moniteur série & Via OpenOCD & STLink uniquement & Via OpenOCD \\
\hline
Licence & Gratuite & Gratuite & Gratuite & Payante \\
\hline
\rowcolor{green!15}
\textbf{Verdict} & \textbf{Retenu} & Non retenu & Non retenu & Non retenu \\
\hline
\end{tabular}
\end{table}

VS~Code associé au framework ESP-IDF en ligne de commande a été
retenu face à Eclipse pour sa légèreté, sa simplicité de configuration
et l'absence de surcharge logicielle sur des projets de grande taille
comme \textit{GlobalFleet} (23 composants ESP-IDF). STM32CubeIDE a été
écarté car il est exclusivement dédié aux microcontrôleurs STM32 et
ne supporte pas la chaîne de compilation RISC-V ESP-IDF.
CLion, bien qu'excellent pour les projets CMake, présente une licence
commerciale payante non justifiée lorsque VS~Code offre des
fonctionnalités équivalentes gratuitement.

%============================================================
\section{Tests et durcissement des firmwares embarqués}
\label{sec:tests}
%============================================================

\subsection{Tests unitaires}
\label{subsec:unit_tests}

Les tests unitaires d'un firmware embarqué constituent le niveau de
vérification le plus granulaire : chaque module ou fonction est testé
de façon isolée, en substituant ses dépendances matérielles par des
implémentations fictives (\textit{mocks}) simulant les comportements
des périphériques. Cette approche permet d'exécuter les tests sur
une machine de développement standard sans matériel cible,
accélérant considérablement le cycle de développement.

Dans ESP-IDF, le framework \textbf{Unity} est intégré nativement
pour l'écriture et l'exécution de tests unitaires \cite{unity_testing}.
Les dépendances matérielles sont remplacées par des mocks implémentant
les mêmes interfaces HAL que les drivers réels. La couche HAL du
firmware \textit{GlobalFleet} a été précisément conçue pour faciliter
cette substitution.

\medskip
\noindent
\textit Note concernant les tests matériels : les tests unitaires
hors-cible du répertoire \texttt{test\_host/} et les tests
d'intégration du répertoire \texttt{test/} ont été entièrement
développés et validés au niveau logiciel. Les tests nécessitant
l'accès au matériel physique ESP32-P4 (profiling de consommation
réelle, validation JTAG, tests de débit CAN physique) n'ont pas pu
être exécutés dans le cadre de ce stage car la carte de développement
matérielle est en cours de finalisation par l'équipe Hardware de
Smart Automation Technologies. Le code de test est prêt et sera
exécuté dès la mise à disposition du matériel. Les résultats seront
intégrés au rapport de la release candidate v1.0.

\subsection{Tests d'intégration}
\label{subsec:integration_tests}

Les tests d'intégration vérifient le comportement correct du système
lorsque plusieurs modules sont combinés, en validant notamment les
interactions via l'Event Bus, les échanges de données entre les files
de messages inter-tâches et le comportement correct des séquences
de démarrage et d'arrêt ordonnés. Contrairement aux tests unitaires
qui s'exécutent sur machine hôte, les tests d'intégration s'exécutent
sur la cible réelle ou sur un émulateur, car ils impliquent les
mécanismes FreeRTOS et les timings réels de l'ordonnanceur.

\subsection{Tests de reboot et de configuration corrompue}
\label{subsec:reboot_tests}

Les tests de robustesse aux coupures d'alimentation et à la corruption
de données consistent à interrompre le système pendant des opérations
d'écriture en Flash, puis à vérifier que le firmware redémarre dans
un état cohérent grâce aux mécanismes de vérification CRC et de
configuration de repli.

Des tests de corruption intentionnelle de la partition NVS via
\texttt{esptool.py} permettent de valider le comportement du module
de configuration en présence de données invalides : le firmware doit
détecter la corruption, journaliser l'incident, basculer sur la
configuration de repli et démarrer normalement sans bloquer ni entrer
dans une boucle de reset infinie.

\subsection{Profiling RAM/stack}
\label{subsec:profiling}

Le profiling de l'utilisation mémoire est une activité de validation
indispensable avant tout déploiement en production. Pour les piles
des tâches FreeRTOS, la fonction
\texttt{uxTaskGetStackHighWaterMark()} retourne le nombre minimum
de mots de pile jamais utilisés simultanément depuis la création de
la tâche. Si ce watermark minimum est inférieur à une marge de
sécurité définie (typiquement 30\,\% de la taille de pile allouée),
la taille est augmentée. Pour le heap, la fonction
\texttt{heap\_caps\_get\_minimum\_free\_size()} mesure le minimum
historique de mémoire libre.

\subsection{Release candidate et documentation}
\label{subsec:release}

La qualification d'une version firmware en tant que release candidate
implique : zéro avertissement de compilation en mode strict
\texttt{-Wall -Wextra}, 100\,\% des tests unitaires en succès,
heap libre supérieur à 100\,Ko au régime pour toutes les variantes,
marge de stack supérieure à 30\,\% pour toutes les tâches, et
signature du binaire avec la clé privée de production. La
documentation technique associée à chaque release candidate comprend
les notes de version détaillant les changements, le rapport de tests
avec les résultats de validation, et la procédure de mise à jour OTA
\cite{unity_testing}.

%============================================================
\section{Conclusion}
\label{sec:conclu_ch2}
%============================================================

Ce chapitre a posé les fondements théoriques et techniques du projet,
en progressant des généralités sur le firmware vers la plateforme
cible, puis vers les mécanismes temps réel et l'architecture logicielle.

L'analyse comparative a justifié trois choix structurants : l'ESP32-P4
(coprocesseur LP, TWAI CAN, OTA Dual-Bank, écosystème ESP-IDF mature),
FreeRTOS (SMP natif, faible empreinte, maturité industrielle) et
Visual Studio Code avec ESP-IDF .

La valeur ajoutée du firmware \textit{GlobalFleet} repose sur trois
contributions : un \textbf{registre de capacités 32 bits} persisté en
NVS, couvrant trois variantes matérielles avec un seul binaire ; un
\textbf{Event Bus Publish-Subscribe} garantissant un couplage nul entre
modules ; et une \textbf{chaîne de récupération graduée} à quatre
niveaux (WARNING → ERROR → CRITICAL → FATAL) assurant la continuité de
service en déploiement distant.

Ces principes — séparation en couches, HAL par inversion de
dépendances, interface uniforme des modules, Event Bus pub/sub, patron
Stratégie — forment la colonne vertébrale de l'architecture détaillée
au chapitre~III.